\documentclass[11pt]{article}

\usepackage[margin=1in]{geometry}
\usepackage[T1]{fontenc}
\usepackage[utf8]{inputenc}
\usepackage{lmodern}
\usepackage{amsmath}
\usepackage{amssymb}
\usepackage{booktabs}
\usepackage{array}
\usepackage{enumitem}
\usepackage{xcolor}
\usepackage{hyperref}

\emergencystretch=2em
\Urlmuskip=0mu plus 2mu

\hypersetup{
  colorlinks=true,
  linkcolor=blue!50!black,
  citecolor=blue!50!black,
  urlcolor=blue!50!black,
  pdftitle={Fensalir: A Fractal Domain Reservoir Pipeline for Resolution-Decoupled Field Rendering},
  pdfauthor={GameCult / Fensalir}
}

\newcommand{\Fensalir}{\textsc{Fensalir}}
\newcommand{\Aquarium}{\textsc{Aquarium}}
\newcommand{\Mimir}{\textsc{Mimir}}
\newcommand{\ReSTIR}{ReSTIR}
\newcommand{\AreaReSTIR}{Area ReSTIR}
\newcommand{\TSR}{TSR}
\newcommand{\GRIS}{GRIS}
\newcommand{\Domain}{\mathcal{D}}
\newcommand{\Sample}{\mathcal{S}}
\newcommand{\Target}{\hat{p}}
\newcommand{\Pdf}{p}
\newcommand{\WeightSum}{W}

\title{\Fensalir: A Fractal Domain Reservoir Pipeline for Resolution-Decoupled Field Rendering}
\author{GameCult / Fensalir}
\date{Living Draft, May 2026}

\begin{document}
\maketitle

\begin{abstract}
We present \Fensalir, a real-time rendering architecture that treats geometry,
material, transport, and sensor observations as typed spatial evidence rather
than as independent pixel-producing passes.  \Fensalir{} combines ideas from
resampled importance sampling, temporal super-resolution, anisotropic splat
rendering, hierarchical virtualized geometry, implicit-field rendering, and
sensor fusion into one pipeline centered on a \emph{fractal domain reservoir}.
Unlike screen-space temporal antialiasing pipelines, the reservoir is not
defined by output texels.  It stores selected samples of semantic field and
radiance claims in their native domains: cube-sphere tiles, object frames,
surface charts, volume domains, tube/curve supports, sensor rigs, and dynamic
hierarchy nodes.  Pixel rows are only the current GPU execution projection of
this deeper sample state.

The core claim is that real-time field rendering should be resolution-decoupled
at the level of sampling authority.  Frame time and resolution may change the
number of probes, candidate samples, hierarchy depth, reuse radius, and
reconstruction aggressiveness, but they must not change which subsystem owns
identity, support, visibility, contribution mass, or temporal validity.  This
paper describes \Fensalir's pipeline, data model, reservoir update equations,
domain-shift validation rules, reconstruction contract, resource-budget model,
and implementation strategy.  It is written as a living research document: the
architecture is confident, the measurements are intentionally explicit about
what has and has not yet been proven, and every named mechanism is tied to the
research archive that shaped it.
\end{abstract}

\section{Introduction}

Real-time rendering pipelines usually inherit a hidden premise: the frame is a
grid of pixels, and temporal systems exist to make that grid look less
undersampled.  Temporal antialiasing accumulates subpixel jitter over frames.
Temporal super-resolution uses richer rejection and reconstruction machinery to
make fewer samples look like more.  Reservoir-based path sampling reuses
selected samples across time and space.  These are powerful ideas, but the
pixel grid remains a gravitational center.

\Fensalir{} starts elsewhere.  It is the native runtime and renderer for
\Aquarium-style clients: authored field grammars, live sensor observations,
direct SDF and tube fields, mesh summaries, density/extinction volumes, and
planetary-scale flow organs all need to become one coherent visible field.  A
pixel is not a natural owner for that evidence.  A spectrum tube has a rolling
buffer coordinate and curve support.  A terrain ridge has a cube-face tile,
grammar path, brush envelope, and conservative summary.  A sensor feature has a
calibration frame, timestamp, modality confidence, and fusion window.  A flame
has density, extinction, emission, and transport uncertainty.  The common
object is not a pixel.  The common object is a bounded claim in a domain.

The thesis of this paper is therefore:

\begin{quote}
\Fensalir{} renders by lowering semantic field and radiance claims into a
hierarchical, budgeted, domain-native reservoir.  Screen-space rows are a
projection of selected evidence, not the foundation of sampling authority.
\end{quote}

This reorients familiar techniques.  \ReSTIR{} showed that one representative
sample can carry the weight mass of many candidates and be reused across
neighboring pixels and frames \cite{bitterli2020restir,ouyang2021restirgi}.
\GRIS{} made the domain and shift-mapping structure explicit
\cite{lin2022gris}.  \AreaReSTIR{} demonstrated that reservoir samples over a
pixel/lens/filter area must store their actual sample coordinates and support
\cite{zhang2024arearestir,areaRestirCode}.  \TSR{}-style reconstruction showed
that temporal history must be inspected, rejected, and reconstructed at the
layer that owns visible output \cite{unrealTSR,unrealTSRFAQ}.  \Fensalir{}
takes the next step: the sample domain is not merely pixel area or lens area,
but the fractal spatial domain tree of the scene itself.

\section{Contributions}

This living draft makes five architectural contributions.

\begin{enumerate}[leftmargin=1.5em]
  \item \textbf{Fractal domain reservoir.}  We define a reservoir sample as a
  selected domain claim carrying identity, selected coordinates, support,
  target value, source proposal policy, represented candidate count,
  contribution weight, conservative bounds, and legal domain shifts.

  \item \textbf{Resolution-decoupled sampling authority.}  We separate the
  native sampling domain from the output grid.  Resolution and frame-time
  targets adjust probe budgets and reconstruction policy, but not semantic
  ownership.

  \item \textbf{Unified field evidence pipeline.}  We place authored IFS/CSG
  claims, direct SDF proxies, mesh summaries, tube fields, density/extinction
  volumes, sensor features, and planetary flow fields behind one bounded
  Form/Appearance/Transport evidence contract.

  \item \textbf{TSR/Area-ReSTIR synthesis.}  We combine domain-coordinate-aware
  reservoir reuse with temporal rejection and reconstruction signals, so the
  denoiser consumes reservoir confidence, support, variance, age, rejection
  reason, and selected sample coordinates rather than opaque pixel history.

  \item \textbf{Budgeted hierarchy and residency model.}  We describe a dynamic
  contribution tree that rebalances node probes with bounded stochastic updates
  and maps selected evidence to CPU, RAM, SSD, and GPU responsibilities.
\end{enumerate}

\section{Prior Work and Research Spine}

\paragraph{Reservoir resampling.}
Resampled importance sampling and weighted reservoir sampling underlie
\ReSTIR's direct-light reuse \cite{bitterli2020restir}.  \ReSTIR{} GI applies
the same spatiotemporal sample reuse shape to indirect paths
\cite{ouyang2021restirgi}.  \GRIS{} generalizes the theory to correlated
samples, domain changes, unknown PDFs, and shift mappings \cite{lin2022gris}.
RTXDI's integration documentation provides practical pass decomposition for
initial sampling, temporal reuse, spatial reuse, and final shading
\cite{rtxdiRestirGI,rtxdiRestirPT}.  \Fensalir{} borrows the reservoir spine
but moves it earlier: primary field candidates themselves are sampled and
reused, not only lighting candidates attached to a known surface.

\paragraph{Area-domain reuse.}
\AreaReSTIR{} extends reservoir reuse to the area of the pixel filter and lens
aperture, storing selected subpixel and lens coordinates and validating reuse
by support overlap, shift mappings, visibility, PDFs, and Jacobians
\cite{zhang2024arearestir,areaRestirCode}.  This is the decisive lesson for
\Fensalir: a tube, SDF splat, density envelope, or fractal tile sample must
carry selected coordinates and support in its own domain.  A reservoir that only
knows ``which texel wrote me'' cannot prove legal reuse.

\paragraph{Temporal reconstruction.}
Modern TAA and \TSR{} systems demonstrate that temporal history needs explicit
reprojection, disocclusion classification, clamping/rejection, accumulated
sample accounting, and debug visualization \cite{unrealTSR,unrealTSRFAQ}.
\Fensalir{} treats these as reservoir resolve responsibilities.  A downstream
presentation pass may filter the resolved field, but it must not become a
second owner of producer identity or temporal persistence.

\paragraph{Implicit fields and splats.}
Hart's sphere tracing establishes the safety contract for implicit surfaces:
distance steps require conservative bounds \cite{hart1996sphere}.  EWA surface
and volume splatting show how anisotropic kernels and screen-space filtering
avoid aliasing and excessive blur \cite{zwicker2001ewa,ren2002objectewa}.
3D Gaussian Splatting and Octree-GS demonstrate the practical value of
anisotropic primitives, visibility sorting, and hierarchy-aware splat LOD
\cite{kerbl2023gaussian,ren2024octreegs}.  \Fensalir{} adopts compact-support
anisotropic envelopes and projected-support scoring, while rejecting infinite
Gaussian tails as a default hot-path cost.

\paragraph{Hierarchical geometry and residency.}
Geometry clipmaps keep terrain detail in nested windows around the viewer
\cite{losasso2004geometry,asirvatham2005gpuclipmaps}.  Nanite demonstrates
fine-grained hierarchical cuts driven by projected error and resident summaries
\cite{karis2021nanite}.  Instant-NGP shows that explicit multiresolution data
structures can make stochastic spatial learning fast and GPU-friendly
\cite{mueller2022instantngp}.  \Fensalir{} combines these into an occupancy and
contribution graph: a hierarchy of conservative summaries, online contribution
estimates, residency state, and selected cuts.

\paragraph{Domain addressing.}
Planetary fields use cube-sphere or quadrilateralized spherical cube domains to
avoid latitude-longitude singularities \cite{projQSC,nasaQSC,clarberg2018cube}.
The Fensalir archive currently favors tangent cube-sphere projection as a
measured first default, with QSC/COBE-like mappings retained for higher
area-regularity paths \cite{fensalirProjectionReport}.  This is not merely a
planet trick.  It is the canonical example of a domain-first renderer: the
surface chart owns the address before the renderer owns packets.

\paragraph{Sensor and flow fields.}
Curless and Levoy's volumetric fusion and KinectFusion show that streamed
sensor observations can accumulate into implicit fields
\cite{curless1996volumetric,newcombe2011kinectfusion}.  Fensalir's planetary
flow organ extends the same evidence vocabulary to stochastic advection fields:
cube-sphere domains, dataset provenance, transition kernels, uncertainty, and
runtime packets \cite{fensalirPlanetaryFlow}.  Its archive cites PlasticAdrift,
GEBCO, ETOPO, ERA5, GLORYS, OSCAR, HYCOM, GraphCast, SFNO/FourCastNet,
MeshGraphNets, GenCast, and NeuralGCM as external anchors for data and model
choices \cite{plasticadrift,mcadam2018transition,gebco2025,etopo2022,era5,
glorys,oscar,hycom,graphcast,sfno,fourcastnet,meshgraphnets,gencast,neuralgcm}.

\paragraph{Presentation and implementation substrate.}
The renderer remains scene-linear through exposure, bloom, and display
transform, following ACES/filmic/HDR production guidance
\cite{aces,epicAutoExposure,unrealBloom,jimenez2014,hableFilmic,
hablePiecewise,hableUncharted,lagarde2014frostbite}.  D3D12 implementation
follows explicit descriptor, barrier, memory, fence, and performance guidance
\cite{d3d12ResourceBinding,d3d12Fences,d3d12Memory,d3d12Interfaces,
directxBarriers,nvidiaBarriers,amdRdna,amdDma,amdGpuUpload}.  These are not the
novel parts of the pipeline, but they constrain the machine: no research
architecture survives contact with a GPU if resource ownership is fog.

\section{System Overview}

\Fensalir{} consumes producer intent and emits a resolved HDR field texture plus
diagnostics.  Producers may be authored DSL grammars, direct SDF shaders,
generated TubeField resources, mesh packages, volume textures, surface pages,
Mimir sensor observations, or planetary simulation outputs.  They do not draw
final truth directly.  They publish bounded claims.

\begin{center}
\begin{minipage}{0.92\linewidth}
\small
\begin{verbatim}
DSL / producer intent / resident resource / live sensor input
-> domain binding
-> semantic Form / Appearance / Transport claims
-> ownership tree and conservative summaries
-> stochastic contribution cache and residency scheduler
-> bounded proposals with target and sourcePdf
-> fractal domain reservoir
-> temporal domain validation and shift
-> spatial/domain validation and shift
-> TSR-style rejection, reconstruction, and denoiser features
-> backend packets and render passes
-> resolved HDR field texture, debug telemetry, evidence ledger
\end{verbatim}
\end{minipage}
\end{center}

The important ownership rule is simple: a backend is a lowering, not a truth
owner.  A mesh may be a first-class contributor if it provides stable identity,
motion or previous-frame mapping, conservative bounds, material/field encoding,
and guide data.  A mesh that only paints pixels is a fallback draw.  The same
rule applies to SDF proxies, TubeField paths, volumes, surface pages, and
sensor fields.

\section{Data Model}

\subsection{Domains}

A domain says where coordinates live and how they map to a parent space:

\begin{center}
\begin{tabular}{ll}
\toprule
Field & Meaning \\
\midrule
\texttt{DomainKey} & stable identity \\
\texttt{ParentKey} & lineage for shift/reuse \\
\texttt{Kind} & cube tile, surface chart, object, volume, tube, sensor rig \\
\texttt{Frame} & local basis and transform \\
\texttt{Projection} & mapping to parent/world coordinates \\
\texttt{Bounds} & conservative support in parent/domain space \\
\texttt{Periodicity} & torus, wrap, seam, or none \\
\texttt{Owner} & engine/client/producible authority \\
\bottomrule
\end{tabular}
\end{center}

Domains are the first validation surface.  If two samples do not have a legal
lineage or shift mapping, they do not reuse each other.

\subsection{Claims}

Claims are semantic statements, not render packets:

\begin{center}
\begin{minipage}{0.86\linewidth}
\small
\begin{verbatim}
Claim {
  ClaimKey, DomainKey, NodeKey,
  Layer: Form | Appearance | Transport,
  Encoding: Height | SignedDistance | Density | Extinction |
            Material | Phase | Emission | Radiance | Feature | Confidence,
  LocalFrame, Envelope, Payload, Tags, CostTier, Seed
}
\end{verbatim}
\end{minipage}
\end{center}

The Form/Appearance/Transport split replaces the earlier implementation
shorthand of separate SDF, PBR, and Radiosity packet names.  Form states what exists where; Appearance
states local material and emission behavior; Transport states how energy moves
through or from the field.  Opaque SDF surfaces, transparent density volumes,
sensor confidence regions, mesh summaries, and spline fields can therefore
share machinery without blending incompatible meanings.

\subsection{Reservoir Samples}

The fractal domain reservoir stores a selected sample:

\begin{center}
\begin{minipage}{0.86\linewidth}
\small
\begin{verbatim}
ReservoirSample {
  SelectedClaim,
  DomainKey,
  SelectedDomainCoordinate,
  SupportFootprint,
  SelectedTarget,
  WeightSum,
  CandidateCount,
  ContributionWeight,
  SourcePdfOrPolicy,
  ProposalKind,
  ShiftMappingMetadata,
  ConservativeBounds,
  Confidence,
  SampleAge,
  RejectionOrValidationMask
}
\end{verbatim}
\end{minipage}
\end{center}

The selected coordinate and support are mandatory for area-domain contributors.
For a TubeField, this includes selected subpixel and curve/buffer-domain
coordinates.  For a cube-sphere terrain sample, it includes face/tile/local UV
and brush support.  For a sensor feature, it includes sensor-rig/world mapping
and calibration confidence.  For a volume, it includes local volume coordinate
and extinction/density support.

\section{Reservoir Update}

\Fensalir{} uses the usual RIS reservoir shape, but the sample type is a field
claim rather than a direct-light sample.  For a candidate \(c\) with target
\(\Target(c)\), source proposal density \(\Pdf(c)\), and represented candidate
count \(m_c\), the unnormalized candidate weight is

\begin{equation}
  w_c = \frac{\Target(c)}{\Pdf(c)}.
\end{equation}

The reservoir maintains

\begin{equation}
  \WeightSum \leftarrow \WeightSum + w_c,\qquad
  M \leftarrow M + m_c,
\end{equation}

and selects the candidate with probability

\begin{equation}
  P(c \text{ selected}) = \frac{w_c}{\WeightSum}.
\end{equation}

The contribution weight exported for final evaluation follows

\begin{equation}
  C = \frac{\WeightSum}{M \cdot \Target(s)},
\end{equation}

where \(s\) is the selected sample.  Invalid zero, negative, infinite, or NaN
targets/PDFs are rejected before proposal emission.  Deterministic structural
proposals use an explicit source policy, commonly \(\Pdf=1\), rather than
pretending to be free samples.

Merging another reservoir \(R_o\) into the current reservoir \(R\) uses the
standard reservoir merge probability:

\begin{equation}
  P(R_o \text{ selected}) =
  \frac{\WeightSum_o}{\WeightSum + \WeightSum_o}.
\end{equation}

The merge is legal only after the candidate reservoir's selected sample has
been shifted or re-evaluated in the target domain and its target/source support
story remains valid.

\section{Domain Shifts and Validation}

Reuse is guilty until proven compatible.  A temporal or spatial pass may not
merge a reservoir because two pixels are nearby, two travels are similar, or two
field ids match.  Those are hints.  Legal reuse requires a domain shift.

\subsection{Temporal Reuse}

Temporal reuse validates:

\begin{itemize}[leftmargin=1.5em]
  \item domain lineage and selected sample support;
  \item camera/object motion and previous-domain reprojection;
  \item disocclusion and occlusion;
  \item travel/depth and normal/gradient compatibility;
  \item field id, material class, layer, and encoding;
  \item motion, previous UV/domain coordinate, and guide confidence;
  \item sample age, variance, and producer-domain validity.
\end{itemize}

The \AreaReSTIR{} lesson is explicit: if the previous selected sample does not
lie in the target support, the reuse pass either rejects it or applies a legal
shift/reconnection with proper evaluation.  The previous reservoir cannot
become visible fallback by itself.

\subsection{Spatial and Domain Reuse}

Spatial reuse is a domain-neighbor problem, not only a screen-neighbor problem.
Valid neighborhoods include screen tiles, cube-face seams, quadtree siblings,
periodic torus wraps, curve supports, volume cells, object-local frames, and
calibrated sensor regions.  For each neighbor, \Fensalir{} evaluates a shift:

\begin{equation}
  T_{a\rightarrow b}: \Sample_a \rightarrow \Sample_b.
\end{equation}

The shifted sample must be re-evaluated in the target domain \(b\).  If the
mapping changes measure, the shift metadata must carry a Jacobian or an
equivalent MIS correction.  If the mapping is deterministic structural replay,
the proposal policy must say so explicitly.

\section{Hierarchy: The Dynamic Contribution Tree}

The reservoir does not float in isolation.  It sits inside an occupancy and
contribution graph:

\begin{center}
\begin{minipage}{0.86\linewidth}
\small
\begin{verbatim}
Node {
  NodeKey, DomainKey, Bounds,
  SummaryPayload, ChildPayload,
  MeanContribution, Variance, Confidence,
  SampleAge, ResidencyState,
  LastVisibleFrame, LastUpdateFrame
}
\end{verbatim}
\end{minipage}
\end{center}

Each node has a conservative summary.  Missing children render through parent
summaries.  The selected hierarchy cut is driven by projected error, material
delta, current reservoir confidence, uncertainty, and estimated CPU/GPU/RAM/SSD
cost.

The projected contribution score begins as a conservative heuristic:

\begin{align}
  r_{\mathrm{px}} &= \mathrm{projectRadius}(B_n, camera),\\
  e_{\mathrm{px}} &= \mathrm{projectError}(E_n, B_n, camera),\\
  S_n &= \frac{
      \max(e_{\mathrm{px}}\cdot V_n,\ \Delta m_n\cdot A_n)
      \cdot I_n
    }{\max(C_n,\epsilon)},
\end{align}

where \(B_n\) is node bound, \(E_n\) is maximum Form error, \(V_n\) is visible
coverage estimate, \(\Delta m_n\) is material delta, \(A_n\) is projected area,
\(I_n\) is importance bias, and \(C_n\) is estimated update/render cost.

\subsection{Bounded Stochastic Probing}

The tree cannot update every node every frame.  \Fensalir{} uses an online
estimator with stochastic exploration:

\begin{equation}
  P_{\mathrm{update}}(n) \propto
  \frac{
    P_{\mathrm{visible}}(n)
    \cdot \mathrm{stale}(n)
    \cdot \mathrm{uncertainty}(n)
    \cdot \max(S_n, S_{\mathrm{parentBias}})
  }{\max(C_n,\epsilon)}.
\end{equation}

On observation \(o\),

\begin{align}
  \delta &= o - \mu_n,\\
  \mu_n &\leftarrow \mu_n + \alpha \delta,\\
  \sigma_n^2 &\leftarrow (1-\beta)\sigma_n^2 + \beta\delta^2.
\end{align}

On non-update, confidence decays and uncertainty grows.  This is deliberately
the boring first machine: EMA, variance, and bandit-style exploration.  A
learned predictor may later rank update and residency priority, but it may not
replace conservative summaries, SDF bounds, calibration bounds, or visibility
safety.

\section{Backend Lowerings}

\Fensalir{} can lower the same evidence contract into multiple execution paths:

\begin{itemize}[leftmargin=1.5em]
  \item 2D height/material/signed-distance/confidence pages;
  \item 3D compact SDF splats for opaque Form;
  \item density/extinction/emission splats for transparent Form and Transport;
  \item 2D fields projected onto 3D surfaces;
  \item direct SDF proxy draws;
  \item generated TubeField curve/tube meshes and analytic tube SDF proposals;
  \item resource-backed mesh packages with stable ids and guide outputs;
  \item sensor confidence volumes and calibrated feature fields;
  \item debug overlays and receipts.
\end{itemize}

The first implementation already has pieces of this: shaped 2D brush envelopes,
semantic claim/tree rows, conservative summaries, fractal contribution state,
projected SDF splat compilation, GPU splat/reservoir receipt harness, direct
TubeField resource lowerings, surface, volume, and mesh resource declarations,
and a shared four-row field reservoir execution ABI.  The live ABI currently
uses:

\begin{center}
\begin{tabular}{cl}
\toprule
Row & Meaning \\
\midrule
0 & current-frame initial RIS \\
1 & temporal reuse reservoir \\
2 & spatial/domain reuse reservoir \\
3 & final resolved reservoir consumed by bloom/presentation \\
\bottomrule
\end{tabular}
\end{center}

Each row is now a nine-lane \texttt{FieldReservoirSample}: color/travel,
metadata, control, guide, motion, domain sample, domain support, proposal, and
stats.  The domain sample lane stores selected pixel/sample UV plus packed
producer coordinates; the domain support lane stores support footprint, domain
kind, and shift kind.  TubeField proposals write selected subpixel UV,
logical source column, curve coordinate, tube support, and explicit motion when
available.  Field id still carries physical rolling-column identity.  SDF
proposals reconstruct object-local hit coordinates and mark object-motion
replay.  Temporal and spatial reuse now reject missing domain state and include
support overlap in the validation weight.  The first TubeField replay path
now re-evaluates the shifted selected logical column/curve coordinate through
a bounded producer-keyed replay manifest and raw source descriptor table.  SDF object-hit
replay reconstructs the previous selected world hit from stored UV/travel,
carries it through object center motion, and validates current ray/travel/support
before temporal reuse can merge it; exact client SDF distance replay remains a
producer-owned extension point.

\section{Reconstruction and Denoising}

\Fensalir{} does not place a conventional TAA pass after rendering and hope it
discovers identity.  Reservoir resolve consumes:

\begin{itemize}[leftmargin=1.5em]
  \item selected domain sample coordinates;
  \item selected support/filter footprint;
  \item contribution weight, weight sum, selected target, and candidate count;
  \item reservoir confidence, sample age, variance, and update probability;
  \item field id, layer, encoding, normal/gradient, travel, and motion;
  \item rejection reason, support miss, disocclusion, occlusion, and shift kind.
\end{itemize}

The denoiser is therefore part of the reservoir pipeline, not downstream
cosmetics.  If a region is rejected because support is missing, it should be
reconstructed differently from a region rejected because of occlusion, field-id
mismatch, low confidence, or high variance.  The scheduler can also respond:
persistent reconstruction stress increases future probe probability in the
relevant domain nodes.

This closes the loop:

\begin{center}
\begin{minipage}{0.86\linewidth}
\small
\begin{verbatim}
reservoir confidence / rejection / variance / denoiser stress
-> contribution cache update
-> future probe and residency priority
-> better selected claims
-> less reconstruction debt
\end{verbatim}
\end{minipage}
\end{center}

\section{Resource Contract}

\begin{table}[h]
\centering
\begin{tabular}{p{0.15\linewidth}p{0.72\linewidth}}
\toprule
Resource & Owned responsibility \\
\midrule
CPU & grammar expansion, domain math, coarse scoring, stochastic scheduling,
selected cuts, page requests, eviction, evidence logs \\
RAM & summaries, estimator state, warm metadata, selected cuts, recent probes \\
SSD & serialized payload pages, summaries, probe history, optional training
datasets; never frame-blocking \\
GPU & selected packet evaluation, SDF/tube/mesh/volume passes, compute scoring
where profitable, page-table sampling, reservoir rows, debug views \\
\bottomrule
\end{tabular}
\caption{The resource contract.  The architecture is budget-adaptive because
each resource has a narrow job.}
\end{table}

Lower frame budgets reduce probes, node depth, candidate count, richer shift
paths, and denoiser feature quality.  They do not change domain ownership.
High budgets buy more evidence.  Low budgets render parent summaries and older
validated reservoirs.  Neither mode lets a pixel row become truth.

\section{Implementation Status}

\Fensalir{} is not a paper-only machine.  The archive records working slices:

\begin{itemize}[leftmargin=1.5em]
  \item cube-sphere projection harness with tangent projection measured against
  normalize projection \cite{fensalirProjectionReport};
  \item shaped compact-support brush envelopes and height-field lowering
  \cite{fensalirFractalRoadmap};
  \item semantic contracts for domains, claims, ownership trees, summaries, and
  selected cuts \cite{fensalirArchitecturePlan};
  \item contribution cache and occupancy graph scaffolding;
  \item Form encodings for signed distance, density, and extinction;
  \item projected 2D-to-3D SDF splat compiler;
  \item GPU receipt harness with two million splats and two million reservoirs
  per layer resident, updating 50,000 entries per layer per frame with two
  candidates per update, measured at 4.471 ms/frame over 120 frames on a GTX
  1070-class machine \cite{fensalirPerfectMachine};
  \item D3D12 TubeField resource and generated-mesh lowerings;
  \item shared field reservoir ABI with current/temporal/spatial/final row
  semantics, selected domain/sample coordinates, support footprints, proposal
  lanes, stats lanes, and row-3 debug modes for rejection, proposal, domain,
  support, and shift state;
  \item HDR presentation path with exposure, bloom pyramid, ACES-fit transform,
  and debug modes \cite{fensalirHDR}.
\end{itemize}

The known architectural gap has moved forward: the current row ABI now carries
domain/sample-coordinate authority, TubeField replay uses a bounded
producer-keyed manifest/source table to re-evaluate selected logical
column/curve coordinates in the current frame, and SDF object-hit replay carries
the previous selected hit through object motion before validating current
ray/travel/support.  The next proof cuts are exact producer SDF replay, measured
temporal leakage/ghosting, and measured rejection precision.
The reservoir work itself is no longer allowed to grow to native presentation
resolution by accident: scene evidence targets, reservoir candidate/history
rows, reservoir resolve, bloom sources, and match-window graph targets are tied
to a clamped internal work grid, while the final present buffer owns full
resolution reconstruction and display.
Those captures now have an explicit comparison switch: native mode enables the
domain validity, support-overlap, and producer-replay gates, while texel-baseline
mode uses the same row ABI but disables those native-domain gates during
temporal reuse.  The first paired Fensalir capture set
\texttt{reservoir-compare-20260529-231904-*} produced a 10.8208\% final-frame
pixel delta.  The repeatable script
\texttt{scripts/measure-reservoir-mode-comparison.ps1} reports that delta; it
is a live-switch check, not a quality claim.
The sequence harness can now prefer a temporal-owner disocclusion mask over the
broad rejection debug view.  That mask is derived from the reservoir rejection
codes for previous history outside the frame, closer previous history, and
support loss; it localizes where an eventual high-budget reference score should
look first.  It can also capture a same-time native-domain reference at a larger
reservoir work-grid scale, producing reference-error rows without introducing a
native-resolution intermediate reservoir.
Those rows feed an offline budget-pressure probe that ranks fixed tiles by
masked reference error, giving the future scheduler a concrete pressure signal
without letting a debug script own runtime visibility.
The renderer now has a fixed spatial-reuse sampling budget in the history update
owner: current and temporal rows remain per-work-grid-pixel, while 3x3 spatial
neighbor reuse runs only on the active budget phase for the frame.
The system has cut the false ``nearest-travel-minus-coverage'' priority owner
and the side-cache temptation.

\section{Evaluation Plan}

The pipeline should be evaluated at four levels.

\paragraph{Reservoir correctness.}
Unit tests should verify RIS add/merge selection, contribution weight,
candidate count preservation, invalid target/PDF rejection, deterministic RNG,
and CPU/GPU parity for packing and update rules.

\paragraph{Domain reuse.}
Fixtures should exercise cube-face seams, quadtree neighbors, tube support,
SDF proxy movement, mesh summaries, density/extinction regions, sensor
calibration shifts, and previous-frame motion.  Negative checks matter: an old
texel path must no longer be able to override the selected domain owner.

\paragraph{Visual behavior.}
Captures should compare final color and debug views for ghosting, occluded-tube
leakage, background speckle, disocclusion rejection, spatial reuse, contribution
weight, support overlap, and denoiser stress.

\paragraph{Budget scaling.}
Benchmarks should sweep resolution and target frame budgets while recording
probe count, selected cut size, GPU time, page requests, stale confidence,
variance, reconstruction debt, and visible error.  The expected result is not
constant quality.  The expected result is graceful degradation without authority
moving back to screen texels.

\section{Discussion}

\Fensalir{} is intentionally vertical.  Its DSL, hierarchy, cache, reservoir,
renderer, temporal reconstruction, denoiser, resource residency, and debug
ledger are designed as one machine.  That is a risk: vertical integration can
become a private cathedral of assumptions.  The countermeasure is explicit
ownership.  Candidate generators own proposal distributions.  Target evaluators
own importance.  Reservoirs own resampling math.  Reuse passes own validation
and shift mappings.  Conservative summaries own safety.  Residency owns
resource pressure.  Resolve owns reconstruction.  Debug surfaces must observe
the same layer where the user sees failure.

This is also why \Fensalir{} does not simply ``use \ReSTIR{}'' or ``use
\TSR{}''.  Those systems solve crucial subproblems, but the renderer's native
objects are broader: domain claims, selected hierarchy nodes, density fields,
sensor features, flow kernels, and packet lowerings.  The proper synthesis is
not another post-process.  It is a domain reservoir with reconstruction-aware
feedback.

\section{Conclusion}

\Fensalir{} proposes a real-time pipeline in which pixels are consumers of
domain evidence, not owners of sampled truth.  A fractal domain reservoir stores
selected field and radiance claims, their support, their contribution mass, and
their legal reuse story.  A dynamic contribution tree decides which domains
deserve more probes under budget.  TSR-style reconstruction consumes the actual
reservoir diagnostics and feeds sampling pressure back into the scheduler.

The result is a renderer that can scale toward arbitrary resolution and frame
targets by adjusting evidence budgets rather than changing the architecture.
The machine is still being built, but the invariant is now sharp enough to cut
against implementation drift: if a sample cannot name its domain, selected
coordinate, support, target, proposal policy, represented mass, and legal
shift, it is not yet a \Fensalir{} reservoir sample.

\clearpage
\begin{thebibliography}{99}
\small

\bibitem{bitterli2020restir}
B. Bitterli, C. Wyman, M. Pharr, P. Shirley, A. Lefohn, and W. Jarosz.
\newblock Spatiotemporal reservoir resampling for real-time ray tracing with dynamic direct lighting.
\newblock \emph{ACM Transactions on Graphics}, 2020.
\newblock \url{https://research.nvidia.com/labs/rtr/publication/bitterli2020spatiotemporal/}

\bibitem{ouyang2021restirgi}
Y. Ouyang, S. Liu, M. Kettunen, M. Pharr, and J. Pantaleoni.
\newblock ReSTIR GI: Path resampling for real-time path tracing.
\newblock \emph{Computer Graphics Forum}, 2021.
\newblock \url{https://research.nvidia.com/publication/2021-06_restir-gi-path-resampling-real-time-path-tracing}

\bibitem{lin2022gris}
D. Lin, M. Kettunen, B. Bitterli, J. Pantaleoni, C. Yuksel, and C. Wyman.
\newblock Generalized resampled importance sampling: Foundations of ReSTIR.
\newblock \emph{ACM Transactions on Graphics}, 2022.
\newblock \url{https://graphics.cs.utah.edu/research/projects/gris/sig22_GRIS.pdf}

\bibitem{zhang2024arearestir}
S. Zhang, D. Lin, M. Kettunen, C. Yuksel, and C. Wyman.
\newblock Area ReSTIR: Resampling for real-time defocus and antialiasing.
\newblock \emph{ACM Transactions on Graphics}, 2024.
\newblock \url{https://graphics.cs.utah.edu/research/projects/area-restir/AreaReSTIR-SIGGRAPH2024.pdf}

\bibitem{areaRestirCode}
S. Zhang.
\newblock Area-ReSTIR source code repository.
\newblock \url{https://github.com/guiqi134/Area-ReSTIR}

\bibitem{rtxdiRestirGI}
NVIDIA.
\newblock RTXDI ReSTIR GI documentation.
\newblock \url{https://github.com/NVIDIA-RTX/RTXDI/blob/main/Doc/RestirGI.md}

\bibitem{rtxdiRestirPT}
NVIDIA.
\newblock RTXDI ReSTIR PT documentation.
\newblock \url{https://github.com/NVIDIA-RTX/RTXDI/blob/main/Doc/RestirPT.md}

\bibitem{unrealTSR}
Epic Games.
\newblock Temporal Super Resolution in Unreal Engine.
\newblock \url{https://dev.epicgames.com/documentation/unreal-engine/temporal-super-resolution-in-unreal-engine}

\bibitem{unrealTSRFAQ}
Epic Games.
\newblock Temporal Super Resolution Frequently Asked Questions.
\newblock \url{https://dev.epicgames.com/documentation/unreal-engine/temporal-super-resolution-frequently-asked-questions-for-unreal-engine}

\bibitem{hart1996sphere}
J. C. Hart.
\newblock Sphere tracing: A geometric method for the antialiased ray tracing of implicit surfaces.
\newblock \emph{The Visual Computer}, 1996.
\newblock \url{https://experts.illinois.edu/en/publications/sphere-tracing-a-geometric-method-for-the-antialiased-ray-tracing/}

\bibitem{zwicker2001ewa}
M. Zwicker, H. Pfister, J. van Baar, and M. Gross.
\newblock EWA splatting.
\newblock \emph{IEEE Visualization}, 2001.
\newblock \url{https://dash.harvard.edu/bitstream/handle/1/4138240/Zwicker_EWA.pdf}

\bibitem{ren2002objectewa}
L. Ren, H. Pfister, and M. Zwicker.
\newblock Object space EWA surface splatting: A hardware accelerated approach to high quality point rendering.
\newblock \emph{Computer Graphics Forum}, 2002.
\newblock \href{https://publications.ri.cmu.edu/object-space-ewa-surface-splatting-a-hardware-accelerated-approach-to-high-quality-point-rendering/}{CMU publication page}

\bibitem{kerbl2023gaussian}
B. Kerbl, G. Kopanas, T. Leimkuehler, and G. Drettakis.
\newblock 3D Gaussian Splatting for Real-Time Radiance Field Rendering.
\newblock \emph{ACM Transactions on Graphics}, 2023.
\newblock \url{https://arxiv.org/abs/2308.04079}

\bibitem{ren2024octreegs}
K. Ren et al.
\newblock Octree-GS: Towards consistent real-time rendering with LOD-structured 3D Gaussians.
\newblock 2024.
\newblock \url{https://arxiv.org/abs/2403.17898}

\bibitem{losasso2004geometry}
F. Losasso and H. Hoppe.
\newblock Geometry clipmaps: Terrain rendering using nested regular grids.
\newblock \emph{ACM Transactions on Graphics}, 2004.
\newblock \url{https://hhoppe.com/geomclipmap.pdf}

\bibitem{asirvatham2005gpuclipmaps}
A. Asirvatham and H. Hoppe.
\newblock Terrain rendering using GPU-based geometry clipmaps.
\newblock \emph{GPU Gems 2}, 2005.
\newblock \url{https://developer.nvidia.com/gpugems/gpugems2/part-i-geometric-complexity/chapter-2-terrain-rendering-using-gpu-based-geometry}

\bibitem{karis2021nanite}
B. Karis, R. Stubbe, and G. Wihlidal et al.
\newblock A deep dive into Nanite virtualized geometry.
\newblock \emph{Advances in Real-Time Rendering}, SIGGRAPH 2021.
\newblock \url{https://advances.realtimerendering.com/s2021/Karis_Nanite_SIGGRAPH_Advances_2021_final.pdf}

\bibitem{mueller2022instantngp}
T. Mueller, A. Evans, C. Schied, and A. Keller.
\newblock Instant neural graphics primitives with a multiresolution hash encoding.
\newblock \emph{ACM Transactions on Graphics}, 2022.
\newblock \url{https://research.nvidia.com/publication/2022-07_instant-neural-graphics-primitives-multiresolution-hash-encoding}

\bibitem{projQSC}
PROJ contributors.
\newblock Quadrilateralized Spherical Cube.
\newblock \url{https://proj.org/en/stable/operations/projections/qsc.html}

\bibitem{nasaQSC}
F. K. Chan and E. M. O'Neill.
\newblock A quadrilateralized spherical cube earth data base.
\newblock NASA Technical Report, 1975.
\newblock \url{https://ntrs.nasa.gov/citations/19810002572}

\bibitem{clarberg2018cube}
P. Clarberg.
\newblock Cube-to-sphere projections for procedural texturing and beyond.
\newblock \emph{Journal of Computer Graphics Techniques}, 2018.
\newblock \url{https://jcgt.org/published/0007/02/01/paper.pdf}

\bibitem{curless1996volumetric}
B. Curless and M. Levoy.
\newblock A volumetric method for building complex models from range images.
\newblock \emph{SIGGRAPH}, 1996.
\newblock \url{https://graphics.stanford.edu/papers/volrange/volrange.pdf}

\bibitem{newcombe2011kinectfusion}
R. A. Newcombe, S. Izadi, O. Hilliges, D. Molyneaux, D. Kim, A. J. Davison, P. Kohli, J. Shotton, S. Hodges, and A. Fitzgibbon.
\newblock KinectFusion: Real-time dense surface mapping and tracking.
\newblock \emph{ISMAR}, 2011.
\newblock \url{https://www.microsoft.com/en-us/research/publication/kinectfusion-real-time-dense-surface-mapping-tracking/}

\bibitem{aces}
Academy Color Encoding System.
\newblock ACES system overview.
\newblock \url{https://docs.acescentral.com/background/overview/}

\bibitem{epicAutoExposure}
Epic Games.
\newblock How Epic Games is handling auto exposure in 4.25.
\newblock \url{https://www.unrealengine.com/tech-blog/how-epic-games-is-handling-auto-exposure-in-4-25?lang=en-US}

\bibitem{unrealBloom}
Epic Games.
\newblock Bloom in Unreal Engine.
\newblock \url{https://dev.epicgames.com/documentation/en-us/unreal-engine/bloom-in-unreal-engine?application_version=5.6}

\bibitem{jimenez2014}
J. Jimenez.
\newblock Next generation post processing in Call of Duty: Advanced Warfare.
\newblock \url{https://www.iryoku.com/next-generation-post-processing-in-call-of-duty-advanced-warfare/}

\bibitem{hableFilmic}
J. Hable.
\newblock Filmic tonemapping operators.
\newblock \url{https://filmicworlds.com/blog/filmic-tonemapping-operators/}

\bibitem{hablePiecewise}
J. Hable.
\newblock Filmic tonemapping with piecewise power curves.
\newblock \url{https://filmicworlds.com/blog/filmic-tonemapping-with-piecewise-power-curves/}

\bibitem{hableUncharted}
J. Hable.
\newblock Uncharted 2 HDR lighting.
\newblock SIGGRAPH 2010 course material.
\newblock \url{https://advances.realtimerendering.com/s2010/Hable-Uncharted2%28SIGGRAPH%202010%20Advanced%20RealTime%20Rendering%20Course%29.pdf}

\bibitem{lagarde2014frostbite}
S. Lagarde and C. de Rousiers.
\newblock Moving Frostbite to physically based rendering 3.0.
\newblock SIGGRAPH 2014 course material.
\newblock \url{https://cgvr.cs.uni-bremen.de/teaching/cg_literatur/Moving%20Frostbite%20to%20Physically%20Based%20Rendering%203.0%2C%202014%2C%20Sebastien%20Lagarde.pdf}

\bibitem{plasticadrift}
PlasticAdrift.
\newblock PlasticAdrift project.
\newblock \url{https://plasticadrift.org/}

\bibitem{mcadam2018transition}
R. McAdam and E. van Sebille.
\newblock Surface connectivity and transition matrices for ocean plastic transport.
\newblock \url{https://dbc.library.uu.nl/handle/1874/362629}

\bibitem{gebco2025}
GEBCO.
\newblock GEBCO 2025 Grid.
\newblock \url{https://www.gebco.net/data-products-gridded-bathymetry-data/gebco2025-grid}

\bibitem{etopo2022}
NOAA NCEI.
\newblock ETOPO global relief model.
\newblock \url{https://www.ncei.noaa.gov/products/etopo-global-relief-model}

\bibitem{era5}
Copernicus Climate Change Service.
\newblock ERA5 climate reanalysis.
\newblock \url{https://climate.copernicus.eu/climate-reanalysis}

\bibitem{glorys}
Copernicus Marine Service.
\newblock GLORYS global ocean physics reanalysis.
\newblock \url{https://data.marine.copernicus.eu/product-detail/GLOBAL_MULTIYEAR_PHY_001_030}

\bibitem{oscar}
Earth and Space Research.
\newblock OSCAR surface currents.
\newblock \url{https://www.esr.org/data-products/oscar/oscar-surface-currents/}

\bibitem{hycom}
HYCOM Consortium.
\newblock HYCOM data server.
\newblock \url{https://www.hycom.org/}

\bibitem{graphcast}
R. Lam et al.
\newblock Learning skillful medium-range global weather forecasting.
\newblock GraphCast.
\newblock \url{https://deepmind.google/research/publications/22598/}

\bibitem{sfno}
T. Bonev et al.
\newblock Spherical Fourier neural operators: Learning stable dynamics on the sphere.
\newblock \url{https://research.nvidia.com/publication/2023-06_spherical-fourier-neural-operators-learning-stable-dynamics-sphere}

\bibitem{fourcastnet}
A. Pathak et al.
\newblock FourCastNet: A global data-driven high-resolution weather model using adaptive Fourier neural operators.
\newblock \url{https://authors.library.caltech.edu/records/k959a-53q45}

\bibitem{meshgraphnets}
T. Pfaff, M. Fortunato, A. Sanchez-Gonzalez, and P. Battaglia.
\newblock Learning mesh-based simulation with graph networks.
\newblock \url{https://arxiv.org/abs/2010.03409}

\bibitem{gencast}
I. Price et al.
\newblock Probabilistic weather forecasting with machine learning.
\newblock GenCast.
\newblock \url{https://www.nature.com/articles/s41586-024-08252-9}

\bibitem{neuralgcm}
S. Kochkov et al.
\newblock Neural general circulation models for weather and climate.
\newblock \url{https://www.nature.com/articles/s41586-024-07744-y}

\bibitem{d3d12ResourceBinding}
Microsoft.
\newblock Resource binding flow of control in Direct3D 12.
\newblock \url{https://learn.microsoft.com/en-us/windows/win32/direct3d12/resource-binding-flow-of-control}

\bibitem{d3d12Fences}
Microsoft.
\newblock Fence-based resource management.
\newblock \url{https://learn.microsoft.com/en-us/windows/win32/direct3d12/fence-based-resource-management}

\bibitem{d3d12Memory}
Microsoft.
\newblock Memory management in Direct3D 12.
\newblock \url{https://learn.microsoft.com/en-us/windows/win32/direct3d12/memory-management}

\bibitem{d3d12Interfaces}
Microsoft.
\newblock Direct3D 12 interfaces.
\newblock \url{https://learn.microsoft.com/en-us/windows/win32/direct3d12/direct3d-12-interfaces}

\bibitem{directxBarriers}
Microsoft DirectX Developer Blog.
\newblock A look inside D3D12 resource state barriers.
\newblock \url{https://devblogs.microsoft.com/directx/a-look-inside-d3d12-resource-state-barriers/}

\bibitem{nvidiaBarriers}
NVIDIA.
\newblock Advanced API performance: barriers.
\newblock \url{https://developer.nvidia.com/blog/advanced-api-performance-barriers/}

\bibitem{amdRdna}
AMD GPUOpen.
\newblock RDNA performance guide.
\newblock \url{https://gpuopen.com/learn/rdna-performance-guide/}

\bibitem{amdDma}
AMD GPUOpen.
\newblock D3D12 Memory Allocator.
\newblock \url{https://gpuopen.com/d3d12-memory-allocator/}

\bibitem{amdGpuUpload}
AMD GPUOpen.
\newblock Using D3D12 heap type GPU upload.
\newblock \url{https://gpuopen.com/learn/using-d3d12-heap-type-gpu-upload/}

\bibitem{fensalirPerfectMachine}
\Fensalir{} archive.
\newblock Perfect Machine architecture.
\newblock \url{file:///E:/Projects/Fensalir/docs/perfect-machine-architecture.md}

\bibitem{fensalirDomainReservoir}
\Fensalir{} archive.
\newblock Fractal domain reservoir.
\newblock \url{file:///E:/Projects/Fensalir/docs/fractal-domain-reservoir.md}

\bibitem{fensalirEvidenceReservoir}
\Fensalir{} archive.
\newblock Temporal spatial evidence reservoir.
\newblock \url{file:///E:/Projects/Fensalir/docs/temporal-spatial-evidence-reservoir.md}

\bibitem{fensalirArchitecturePlan}
\Fensalir{} archive.
\newblock Fractal brush architecture plan.
\newblock \url{file:///E:/Projects/Fensalir/research/rendering/fractal-brush-architecture-plan.md}

\bibitem{fensalirCubeBrushes}
\Fensalir{} archive.
\newblock Cube-sphere fractal brush research.
\newblock \url{file:///E:/Projects/Fensalir/research/rendering/cube-sphere-fractal-brushes.md}

\bibitem{fensalirFractalRoadmap}
\Fensalir{} archive.
\newblock Fractal brush implementation roadmap.
\newblock \url{file:///E:/Projects/Fensalir/research/rendering/fractal-brush-implementation-roadmap.md}

\bibitem{fensalirProjectionReport}
\Fensalir{} archive.
\newblock Fractal projection report.
\newblock \url{file:///E:/Projects/Fensalir/research/rendering/fractal-projection-report.md}

\bibitem{fensalirHDR}
\Fensalir{} archive.
\newblock HDR/PBR presentation research.
\newblock \url{file:///E:/Projects/Fensalir/research/rendering/hdr-pbr-presentation.md}

\bibitem{fensalirPlanetaryFlow}
\Fensalir{} archive.
\newblock Planetary flow organ literature.
\newblock \url{file:///E:/Projects/Fensalir/research/planetary-flow-organ-literature.md}

\bibitem{fensalirRenderer}
\Fensalir{} archive.
\newblock HLSL renderer documentation.
\newblock \url{file:///E:/Projects/Fensalir/docs/hlsl-renderer.md}

\bibitem{fensalirTSR}
\Fensalir{} archive.
\newblock TSR-inspired TAA specification.
\newblock \url{file:///E:/Projects/Fensalir/docs/tsr-inspired-taa-spec.md}

\bibitem{fensalirTeardown}
\Fensalir{} archive.
\newblock Rendering teardown/rebuild protocol.
\newblock \url{file:///E:/Projects/Fensalir/docs/rendering-teardown-rebuild-protocol.md}

\bibitem{vibeGeometryDoctrine}
VibeGeometry archive.
\newblock Procedural doctrine and CSG research.
\newblock \url{file:///E:/Projects/VibeGeometry/docs/research/procedural-doctrine.md},
\url{file:///E:/Projects/VibeGeometry/docs/research/dream-machine-grammar.md},
\url{file:///E:/Projects/VibeGeometry/docs/research/realtime-csg-doctrine.md}

\bibitem{perfectMachineArticle}
GameCult archive.
\newblock The Perfect Machine Has One Eye.
\newblock \url{file:///E:/Projects/gamecult-site/GameCult/Blog/perfect-machine-fractal-reservoir-architecture.md}

\end{thebibliography}

\end{document}
